\section{Fundamentos del aprendizaje profundo}
\label{sec:marco_fundamentos}

El aprendizaje automático estudia algoritmos que mejoran su desempeño en una tarea a medida que acumulan
experiencia, que en la práctica son datos. En lugar de programar la solución como un conjunto de reglas, se
elige un modelo con parámetros ajustables y se los fija a partir de ejemplos, de modo que el modelo resuelva la
tarea también sobre datos que no vio~\cite{goodfellow2016deep}. En el aprendizaje \textbf{supervisado}, cada
ejemplo es un par formado por una entrada, como una imagen, un texto o un conjunto de mediciones, y la salida
esperada, su etiqueta, y el modelo aprende a predecir la etiqueta de entradas nuevas. En el \textbf{no
supervisado} los ejemplos no tienen etiqueta y se busca la estructura de los propios datos, como grupos de
casos parecidos~\cite{goodfellow2016deep}. El \textbf{autosupervisado} obtiene la etiqueta de los mismos datos,
por ejemplo la palabra que sigue en un texto, y es como se preentrenan los modelos de lenguaje
(Sección~\ref{sec:marco_llm_grandes}).

Durante décadas, el desempeño de estos métodos dependió de las características con que se describe cada
ejemplo. Un especialista las diseñaba a mano, como medidas de color o textura para una imagen o recuentos de
palabras para un texto, y el algoritmo aprendía solo a decidir a partir de ellas. El aprendizaje profundo
aprende también esas características, directamente de los píxeles o de las palabras. Para eso encadena muchas
transformaciones simples, de modo que las primeras capturan rasgos elementales y las últimas, conceptos más
abstractos; en una imagen, de bordes a contornos y de contornos a partes de objetos. El enfoque se impuso a
partir de la década de 2010 en el reconocimiento de imágenes y de voz y en el procesamiento del lenguaje,
cuando hubo suficientes datos y cómputo para entrenar esos modelos~\cite{goodfellow2016deep}, y es la base de
todos los modelos que usa este trabajo. Las transformaciones que encadena son las capas de una red neuronal.

\subsection{Redes neuronales y aprendizaje supervisado}
\label{sec:marco_redes}

Una red neuronal artificial es un modelo formado por unidades de cálculo simples, las neuronas, conectadas
entre sí y organizadas en capas. Cada neurona calcula una suma ponderada de sus entradas, le suma un sesgo y
aplica al resultado una función no lineal, llamada función de activación. Una capa reúne varias neuronas que reciben la misma entrada,
y un perceptrón multicapa apila capas de modo que la salida de cada una es la entrada de la siguiente. La capa
$l$ transforma el vector $\mathbf{h}^{(l-1)}$ que recibe en
\begin{equation}
    \mathbf{h}^{(l)} = \phi\!\left(W^{(l)} \mathbf{h}^{(l-1)} + \mathbf{b}^{(l)}\right),
    \label{eq:marco_capa}
\end{equation}
donde $W^{(l)}$ es la matriz de pesos, $\mathbf{b}^{(l)}$ el vector de sesgos y $\phi$ la función de activación,
aplicada por elemento. Sin la activación, la red entera se reduciría a una sola transformación lineal. Los pesos
y los sesgos de todas las capas son los parámetros de la red, agrupados en $\theta$, y una red de muchas capas
se llama profunda~\cite{goodfellow2016deep}.

Entrenar la red es elegir los valores de $\theta$. Para eso se define una función de pérdida $\mathcal{L}$, que
cuantifica cuánto se equivoca la red en un ejemplo, y se buscan los parámetros que minimizan la pérdida promedio
sobre el conjunto de entrenamiento. El descenso del gradiente lo hace de forma iterativa. El gradiente de la
pérdida respecto de los parámetros, $\nabla_{\theta} \mathcal{L}$, indica en qué dirección crece la pérdida, y
en cada paso los parámetros se desplazan en la dirección opuesta una distancia que fija la tasa de aprendizaje
$\eta$. La regla de actualización es entonces
\begin{equation}
    \theta \leftarrow \theta - \eta \, \nabla_{\theta} \mathcal{L}.
    \label{eq:marco_gradiente}
\end{equation}
Cada paso estima el gradiente sobre un lote de ejemplos y no sobre el conjunto completo, y una época es una
pasada por todos los lotes~\cite{goodfellow2016deep}.

El gradiente de todos los parámetros se obtiene con la retropropagación, que recorre la red desde la pérdida
hacia la entrada aplicando la regla de la cadena (Figura~\ref{fig:marco_perceptron}), con un costo del mismo orden que el de calcular la
salida~\cite{rumelhart1986learning}. Sobre la regla de la Ecuación~\ref{eq:marco_gradiente} se construyen
optimizadores como Adam, que adapta el paso de cada parámetro con dos promedios móviles, uno del gradiente y
otro de su cuadrado~\cite{kingma2015adam}. Esos promedios son memoria adicional por cada parámetro que se
entrena, un costo que pesa al adaptar modelos grandes (Sección~\ref{sec:marco_llm_lora}). La tasa de
aprendizaje suele crecer de forma lineal en las primeras épocas, un calentamiento (\emph{warmup}) que evita
pasos grandes mientras la red cambia rápido~\cite{goyal2017accurate}, y decaer después hasta cerca de cero
siguiendo medio ciclo de coseno~\cite{loshchilov2017sgdr}. Una vez entrenada, la red predice la salida de una
entrada nueva con una pasada hacia adelante y los parámetros fijos, lo que se llama inferencia.

% Figura: perceptrón multicapa, pasada hacia adelante y retropropagación.
% Se usa en sec:marco_fundamentos, subsección de entrenamiento.

\begin{figure}[H]
    \centering
    \begin{tikzpicture}[x=1mm, y=1mm,
        neurona/.style = {circle, draw=figtrazo, minimum size=6mm, inner sep=0pt, font=\footnotesize},
    ]

        % Capas.
        \foreach \i in {1, 2, 3} {
            \node[neurona, fill=figazul!15, draw=figazul!80!black] (x\i) at (0, 16 - \i*8) {$x_{\i}$};
        }
        \foreach \i in {1, 2, 3, 4} {
            \node[neurona, fill=fignaranja!18, draw=fignaranja!70!black] (h\i) at (25, 20 - \i*8) {};
        }
        \foreach \i in {1, 2, 3} {
            \node[neurona, fill=figverde!12, draw=figverde!70!black] (z\i) at (50, 16 - \i*8) {$z_{\i}$};
        }

        % Conexiones.
        \foreach \i in {1, 2, 3} { \foreach \j in {1, 2, 3, 4} {
            \draw[figtrazo!60, thin] (x\i) -- (h\j);
        } }
        \foreach \j in {1, 2, 3, 4} { \foreach \k in {1, 2, 3} {
            \draw[figtrazo!60, thin] (h\j) -- (z\k);
        } }

        % Softmax y pérdida.
        \node[figtarea, text width=16mm, minimum height=8mm] (sm) at (72, 0) {softmax};
        \node[figsalida, text width=20mm, minimum height=8mm] (loss) at (100, 0)
            {pérdida\\$\mathcal{L} = -\log p_y$};
        \draw[figflecha] (56, 0) -- (sm.west);
        \draw[figflecha] (sm.east) -- node[figetiqueta, above] {$p_c$} (loss.west);

        % Rótulos de capas.
        \node[font=\footnotesize, text=figtrazo] at (0, -22) {entrada};
        \node[font=\footnotesize, text=figtrazo, align=center] at (25, -22) {capa oculta\\$\phi(W\mathbf{x}+\mathbf{b})$};
        \node[font=\footnotesize, text=figtrazo] at (50, -22) {\emph{logits}};

        % Pasadas.
        \draw[figflecha, figazul!80!black, thick] (0, 24) -- node[figetiqueta, above, text=figazul!80!black]
            {pasada hacia adelante: se calcula la salida} (100, 24);
        \draw[-{Latex[length=2mm, width=1.6mm]}, thick, densely dashed, figbermellon]
            (100, -32) -- node[figetiqueta, below, text=figbermellon]
            {retropropagación: se calcula $\nabla_\theta \mathcal{L}$ capa por capa} (0, -32);

    \end{tikzpicture}
    \caption{Perceptrón multicapa con una capa oculta: pasada hacia adelante y retropropagación.}
    \label{fig:marco_perceptron}
\end{figure}


\subsection{Clasificación y funciones de pérdida}
\label{sec:marco_softmax}

Un problema de clasificación consiste en asignar cada entrada a una clase de un conjunto finito. En el diagnóstico de lesiones de piel, las clases son los tipos de lesión, como el
melanoma o el nevo, y el modelo recibe la imagen y debe indicar a cuál pertenece. Para un problema de $C$
clases, la última capa de la red tiene una salida por clase y devuelve un vector de $C$ puntajes reales sin
restricción, los \emph{logits} $\mathbf{z}$. La función \emph{softmax} los convierte en una distribución de
probabilidad,
\begin{equation}
    p_c = \frac{\exp(z_c)}{\sum_{j=1}^{C} \exp(z_j)}, \qquad c = 1, \dots, C,
    \label{eq:marco_softmax}
\end{equation}
con valores positivos que suman uno. La clase predicha es la de mayor probabilidad, y esa probabilidad se lee
como la confianza del modelo en su predicción, una lectura que solo vale si el modelo está calibrado
(Sección~\ref{sec:marco_senales_confianza}). Los modelos de lenguaje usan la misma operación para elegir cada
palabra entre las de su vocabulario (Sección~\ref{sec:marco_llm_grandes}).

La pérdida de clasificación es la entropía cruzada. Compara la distribución que predice el modelo,
$\mathbf{p}$, con la de la etiqueta, $\mathbf{q}$, que pone toda la probabilidad en la clase correcta $y$, y
como $q_c$ vale cero para las demás clases, se reduce a menos el logaritmo de la probabilidad asignada a la
correcta,
\begin{equation}
    \mathcal{L} = -\sum_{c=1}^{C} q_c \log p_c = -\log p_y.
    \label{eq:marco_entropia}
\end{equation}
Vale cero cuando toda la probabilidad está en la clase correcta y crece sin límite a medida que esa
probabilidad se acerca a cero, de modo que un error cometido con mucha confianza pesa más que uno cometido con
duda. Minimizarla equivale a maximizar la probabilidad que el modelo asigna a las respuestas
correctas~\cite{goodfellow2016deep}. Combinada con la softmax, su gradiente respecto de cada logit $z_c$ es
$p_c - q_c$: el entrenamiento sube el logit de la clase correcta y baja el de las demás en proporción a la
probabilidad que el modelo les había asignado. Los modelos de lenguaje usan la misma pérdida para cada token
(Sección~\ref{sec:marco_llm_grandes}), y las correcciones para el desbalance de clases
(Sección~\ref{sec:marco_desbalance}) son variantes de ella.

\subsection{Generalización y desafíos del aprendizaje}
\label{sec:marco_regularizacion}

Lo que interesa de un modelo es su desempeño sobre casos que no vio, y una red con suficientes parámetros puede
memorizar el conjunto de entrenamiento: la pérdida de entrenamiento sigue bajando mientras el error sobre datos
nuevos sube, fenómeno que se llama \textbf{sobreajuste}~\cite{goodfellow2016deep}.

Ese error se estima con ejemplos que el modelo no usó para ajustarse, y por eso los datos se dividen en tres
particiones disjuntas. La de entrenamiento ajusta los parámetros, la de validación elige entre configuraciones
e hiperparámetros, y la de test se usa una sola vez al final, porque una decisión tomada mirando el test vuelve
optimista su estimación~\cite{goodfellow2016deep}. Si imágenes de un mismo paciente o de una misma lesión caen
en el entrenamiento y en el test, el modelo puede acertar las segundas por haber visto las primeras y no porque
generalice. Esta fuga de información (\emph{data leakage}) hace la estimación más favorable que el desempeño
real, y se evita si todas las observaciones de un mismo caso caen en la misma partición.

La \textbf{regularización} reduce el sobreajuste restringiendo lo que la red puede aprender o ampliando la
variedad de lo que ve. El decaimiento de pesos (\emph{weight decay}) encoge los parámetros hacia cero en cada
paso, en una proporción $\lambda$, lo que con el descenso del gradiente equivale a sumar a la pérdida la
penalización $(\lambda/2)\sum_i \theta_i^2$~\cite{goodfellow2016deep}. Con Adam esa equivalencia no vale, y
AdamW aplica el decaimiento por separado de la actualización del optimizador~\cite{loshchilov2019decoupled}. La
detención temprana (\emph{early stopping}) conserva los parámetros de la época con mejor desempeño en
validación~\cite{goodfellow2016deep}. El suavizado de etiquetas (\emph{label smoothing}) le quita a la clase
correcta de $\mathbf{q}$ una fracción $\varepsilon$ de la probabilidad y la reparte entre todas las clases,
con lo que el modelo deja de empujar $p_y$ hacia uno y suele quedar mejor calibrado
(Sección~\ref{sec:marco_senales_confianza})~\cite{szegedy2016rethinking,muller2019does}.

El aumento de datos aplica a cada imagen, cada vez que entra en un lote, una transformación aleatoria que
preserva su clase, como una rotación, un recorte o un cambio de brillo. RandAugment aplica $N$ operaciones
sorteadas de una lista fija con una magnitud común $M$~\cite{cubuk2020randaugment}, y el borrado aleatorio
(\emph{random erasing}) tapa un rectángulo al azar para simular oclusiones~\cite{zhong2020random}. Cuando el color distingue las clases, como
en dermatoscopía~\cite{kittler2016standardization}, una transformación de color fuerte puede borrar la
información que las define~\cite{shorten2019survey}. El aumento se usa también al predecir: el \gls{tta}
promedia las probabilidades que el modelo asigna a la imagen original y a algunas versiones transformadas de
ella, como su reflejo horizontal~\cite{krizhevsky2012imagenet}.

Otro desafío es el desbalance de clases\label{sec:marco_desbalance}. Cuando unas clases tienen muchos más
ejemplos que otras, minimizar la entropía cruzada promedio favorece a las frecuentes, porque sus ejemplos
dominan la suma de la pérdida, y el efecto crece con el grado de desbalance~\cite{buda2018systematic}. Sean
$n$ la cantidad total de ejemplos de entrenamiento, $n_c$ la de la clase $c$ e $y$ la clase correcta del
ejemplo. Las correcciones más usadas modifican la entropía cruzada de la Ecuación~\ref{eq:marco_entropia} para
que las clases raras pesen más:
\begin{enumerate}
    \item la \textbf{pérdida ponderada}, $-w_y \log p_y$, multiplica la pérdida de cada ejemplo por un peso
    inversamente proporcional a la frecuencia de su clase, $w_y = n/(C\,n_y)$, de modo que cada clase aporta
    en total el mismo peso a la suma;
    \item la \textbf{pérdida focal}, $-(1-p_y)^{\gamma}\log p_y$, agrega un factor que se acerca a cero en los
    ejemplos que el modelo ya clasifica con confianza, con $p_y$ cerca de uno, y concentra así el
    entrenamiento en los difíciles. El parámetro $\gamma \ge 0$ regula la atenuación, y con $\gamma = 0$ se
    recupera la entropía cruzada~\cite{lin2017focal};
    \item el \textbf{\emph{balanced softmax}} multiplica en la softmax de la Ecuación~\ref{eq:marco_softmax}
    la exponencial de cada logit por la cantidad de ejemplos de su clase,
    $p_c = n_c \exp(z_c) / \sum_{j} n_j \exp(z_j)$, lo que equivale a sumar $\log n_c$ al logit $z_c$. Durante
    el entrenamiento, las clases frecuentes reciben así una ventaja que el modelo no necesita aprender en sus
    logits, y al predecir se usa la softmax común, sin ese término~\cite{ren2020balanced}.
\end{enumerate}
La alternativa es entrenar sin corrección y ajustar la regla de decisión, dividiendo la probabilidad de cada
clase por su proporción en el entrenamiento~\cite{buda2018systematic} o eligiendo en validación un umbral para
una clase rara. Ninguna de estas correcciones agrega información sobre las clases raras. Además, la exactitud
global oculta el desbalance, y por eso se usan métricas que promedian por clase
(Sección~\ref{sec:marco_eval_clasif}): la sensibilidad macro, que mide qué fracción de los casos de cada clase
se detecta, y el F1 macro, que además penaliza los falsos positivos.

\subsection{Modelos preentrenados y adaptación}
\label{sec:marco_transferencia}

Con pocos ejemplos etiquetados, como ocurre en la mayoría de los conjuntos de datos médicos, entrenar una red
desde cero favorece el sobreajuste. La transferencia de aprendizaje reutiliza lo aprendido en una tarea de
origen con muchos datos para mejorar el aprendizaje en la tarea de destino~\cite{pan2010survey}. En
reconocimiento visual, el preentrenamiento se hace habitualmente sobre ImageNet, un conjunto de referencia de
imágenes naturales etiquetadas~\cite{russakovsky2015imagenet,raghu2019transfusion}. Las primeras capas de una
red visual aprenden rasgos generales, parecidos a filtros de bordes y de color, que se reutilizan bien en otras
tareas, y las últimas, rasgos específicos de la tarea de origen~\cite{yosinski2014transferable}. En imágenes
médicas el beneficio puede ser menor al esperado, por la distancia con las imágenes
naturales~\cite{raghu2019transfusion}, y esa distancia motiva los modelos preentrenados sobre imágenes del
propio dominio, como DermLIP (Sección~\ref{sec:marco_vlm_clip}).

Para adaptar una red que clasifica se la divide en dos partes. El \emph{backbone} convierte la imagen en un
vector de características, y la cabeza de clasificación (\emph{head}) es la capa final, de la forma de la
Ecuación~\ref{eq:marco_capa}, que convierte ese vector en los logits. La cabeza se reemplaza por una nueva, con
una salida por cada clase de la tarea de destino. La \textbf{extracción de características} congela el
\emph{backbone}, es decir, mantiene fijos sus parámetros, y entrena solo la cabeza. El \textbf{ajuste fino}
(\emph{fine-tuning}) actualiza también el \emph{backbone}, completo o solo sus últimas capas. Entrenar más
parámetros da más libertad para adaptarse a la tarea nueva, pero con pocos datos aumenta el riesgo de
sobreajuste y de olvido catastrófico (\emph{catastrophic forgetting}), es decir, que el modelo pierda parte de
lo que aprendió en el preentrenamiento.

Entre esos dos extremos hay estrategias que adaptan el \emph{backbone} sin deformar lo que aprendió. La primera
entrena en dos etapas: primero solo la cabeza, con el \emph{backbone} congelado, y después la red completa. Si
el ajuste fino empieza con una cabeza inicializada al azar, los gradientes que llegan al \emph{backbone}
mientras la cabeza aprende deforman las características preentrenadas. Con la cabeza ya entrenada, la
secuencia rinde más que cada una de sus dos partes, sobre todo ante datos de una distribución distinta de la
del entrenamiento~\cite{kumar2022finetuning}. La segunda asigna a cada capa su propia tasa de aprendizaje,
decreciente hacia la entrada, porque las capas capturan información de distinto tipo~\cite{howard2018universal}.
Con $L$ capas numeradas desde la entrada y un factor $\xi < 1$, la capa $l$ se entrena con la tasa
\begin{equation}
    \eta_l = \eta \, \xi^{\,L-l}, \qquad l = 1, \dots, L,
    \label{eq:marco_llrd}
\end{equation}
de modo que la cabeza recibe la tasa completa y las primeras capas, que guardan los rasgos más generales,
cambian menos. La técnica se conoce como decaimiento de la tasa de aprendizaje por capa (\emph{layer-wise
learning rate decay}).

Hasta aquí, la red trata su entrada como un vector sin estructura. Las arquitecturas que hicieron práctico el
aprendizaje profundo incorporaron la estructura de cada tipo de dato, y la primera en hacerlo con las imágenes
fue la red convolucional.
